\documentclass{article}
\usepackage{amssymb,amsmath}
\usepackage{graphicx}
\usepackage{enumerate}
\usepackage{amssymb,amsmath,amsthm}
\usepackage{graphicx}
\usepackage{tikz}
\usepackage{pgfplots}
\usepackage[utf8]{inputenc}
\usepackage[brazil]{babel}
\usepackage{enumerate}
\usepackage{color}
\usepackage{mathabx}
\usepackage{calc}
\usepackage{fullpage}

\usepackage{dsfont} % includes bb 1
\newcommand*\1{\mathds{1}}
\usepackage[brazil]{babel}

\newtheorem{theorem}{Teorema}[section]
\newtheorem{corollary}[theorem]{Corolário}
\newtheorem{lemma}[theorem]{Lema}
\newtheorem{proposition}[theorem]{Proposição}
\newtheorem{definition}[theorem]{Definição}
\newtheorem{notation}[theorem]{Notação}

\DeclareMathOperator{\Var}{Var}
\def\d{\mathrm{d}}

\begin{document}

\title{Probabilidade I}
\author{Primeira Prova}

\maketitle

\noindent Observa\c{c}\~oes:
\begin{itemize}
\item A prova terá duração de $3$ horas e meia.
\item Questões deverão ser respondidas de maneira matematicamente rigorosa e clara.
\item Todos os resultados provados em aula (exceto exercícios) poderão ser utilizados, salvo quando a questão pedir que algum destes seja provado. Nesse caso, todos os resultados anteriores poderão ser utilizados.
\item Caso uma questão seja composta de vários itens, cada um poderá ser resolvido independentemente, supondo a validade dos anteriores.
\end{itemize}

\vspace{4mm}
\noindent {\bf 1)} A distribuição de Fréchet é frequentemente utilizada para modelar extremos, como máximos anuais de vazão de rios ou cargas máximas em estruturas.

\begin{enumerate}[\quad a)]
  \item Sejam $W_1$ e $W_2$ variáveis aleatórias independentes com distribuição de Fréchet de parâmetros $(\alpha,s,m)$, cuja função acumulada de distribuição é dada por
\begin{equation}
  F(x)=
  \begin{cases}
    \quad
    \displaystyle
    \exp\bigg\{
    -\left(\frac{x-m}{s}\right)^{-\alpha}
    \bigg\},
    & x>m,\\
    \\
    \quad
    0, & x \le m.
  \end{cases}
\end{equation}

Defina
\[
Y=\max\{W_1, W_2\}.
\]
 Determine a função acumulada de distribuição de $Y$, mostre que $Y$ também possui distribuição de Fréchet e identifique seus parâmetros.
\item Sejam $X_1,X_2,\dots$ variáveis aleatórias iid com função acumulada de distribuição
\begin{equation}
F(x)=
\begin{cases}
0, & x<1,\\
1-x^{-\alpha}, & x\ge 1,
\end{cases}
\end{equation}
onde $\alpha>0$ e defina
  \begin{equation}
    M_n=\max\{X_1,\dots,X_n\}.
  \end{equation}
  e mostre que
  \begin{equation}
    \frac{M_n}{n^{1/\alpha}}
    \overset{d}{\Longrightarrow}
    Y,
  \end{equation}
  onde $Y$ possui distribuição de Fréchet com parâmetros $(\alpha,1,0)$.
\end{enumerate}


\vspace{4mm}
\noindent {\bf 2)} O grafo aleatório de Erdős--Rényi $G(n,p)$ é o grafo com conjunto de vértices
\begin{equation}
\{1,2,\dots,n\},
\end{equation}
em que cada aresta é incluída independentemente das demais com probabilidade $p$.
Um vértice é dito \emph{isolado} se não existe nenhuma aresta incidente a ele.
Seja $I_n$ o número de vértices isolados em $G(n,p)$ e
\begin{enumerate}[\quad a)]
    \item Calcule $\mathbb E[I_n]$.

    \item Mostre que para todo $\varepsilon > 0$, se
    \begin{equation}
    p_n = (1 + \varepsilon) \frac{\log n}{n},
    \end{equation}
    então
    \begin{equation}
    \mathbb P(I_n=0)\to1.
    \end{equation}
  \item Seja \(X\ge0\) uma variável aleatória com variância finita e tal que $\mathbb E[X] > 0$.
Mostre que
\begin{equation}
\mathbb P(X=0)
\le
\frac{\mathrm{Var}(X)}{\mathbb E[X]^2}.
\end{equation}
    \item Mostre que, para todo $\varepsilon > 0$, se
    \begin{equation}
    p_n = (1 - \varepsilon) \frac{\log n}{n},
    \end{equation}
    então $\mathbb P(I_n>0)\to1$.
\end{enumerate}

\vspace{4mm}
\noindent {\bf 3)} [Lei $\{0,1\}$ de Hewitt--Savage]
Sejam \(X_1,X_2,\dots\) variáveis aleatórias iid com valores em um espaço mensurável \(E\).
Um evento \(A\in\sigma(X_1,X_2,\dots)\) é dito \emph{invariante por permutações finitas} se, para toda permutação \(\pi\) de \(\mathbb N\) que fixa todos exceto um número finito de índices,
\[
\mathbf 1_A(X_1,X_2,\dots)
=
\mathbf 1_A(X_{\pi(1)},X_{\pi(2)},\dots)
\quad \text{q.c.}
\]
O objetivo desta questão é mostrar que $\mathbb P(A)\in\{0,1\}$.

\begin{enumerate}[\quad a)]
    \item Seja \(\varepsilon>0\). Mostre que existe \(n\ge1\) e um evento $B\in\sigma(X_1,\dots,X_n)$ tal que
    \[
    \mathbb P(A\triangle B)<\varepsilon,
    \]
    onde \(A\triangle B\) denota a diferença simétrica.

    \item Para \(m>n\), defina \(B^{(m)}\) como o evento obtido à partir de \(B\), trocando as variáveis $X_1,\dots,X_n$ por $X_{m+1},\dots,X_{m+n}$.
    Isto é, se $B=\{(X_1,\dots,X_n)\in C\}$, então $B^{(m)}=\{(X_{m+1},\dots,X_{m+n})\in C\}$.
    Mostre que
    \[
    \mathbb P(B^{(m)})=\mathbb P(B).
    \]

    \item Use a invariância de \(A\) por permutações finitas para mostrar que
    \[
    \mathbb P(A\triangle B^{(m)})<\varepsilon.
    \]

    \item Conclua que
    \[
    |\mathbb P(A\cap A)-\mathbb P(B\cap B^{(m)})|
    \le 2\varepsilon.
  \]
  e que
    \[
    |\mathbb P(A)-\mathbb P(A)^2|
    \le C\varepsilon
    \]
    para alguma constante universal \(C\).
    Como \(\varepsilon>0\) é arbitrário, conclua que $\mathbb P(A)=\mathbb P(A)^2$.
    Portanto, $\mathbb P(A)\in\{0,1\}$.
\end{enumerate}

\end{document}
